\documentclass{article}

% NeurIPS 2026 style — use [final] for camera-ready, omit for anonymous submission
\usepackage[final]{neurips_2026}

\usepackage[utf8]{inputenc}
\usepackage[T1]{fontenc}
\usepackage{hyperref}
\usepackage{url}
\usepackage{booktabs}
\usepackage{amsfonts}
\usepackage{amsmath}
\usepackage{nicefrac}
\usepackage{microtype}
\usepackage{xcolor}
\usepackage{graphicx}
\usepackage{algorithm}
\usepackage{algpseudocode}
\usepackage{multirow}

\title{ALVRO: Adaptive LVR Optimizer ---\\
A Reinforcement Learning Agent for\\
Concentrated Liquidity Provision on Uniswap~v3}

\author{%
  Tong LI \\
  AIvancity School of Technology, Business \& Society \\
  Paris, France \\
  \texttt{cleo.tongli@gmail.com} \\
}

\begin{document}

\maketitle

% ─────────────────────────────────────────────────────────────────────────────
\begin{abstract}
Providing liquidity in decentralised exchanges (DEXs) exposes liquidity
providers (LPs) to \emph{Loss-Versus-Rebalancing} (LVR), a continuous value
leak to arbitrageurs that is unavoidable under passive management.  We
introduce \textbf{ALVRO} (Adaptive LVR Optimizer), a reinforcement learning
agent that actively manages a concentrated liquidity position on Uniswap~v3
to maximise fee income while minimising LVR and gas friction.  ALVRO combines
three components: (i)~a GARCH(1,1) conditional volatility model for real-time
risk estimation; (ii)~a \emph{DeepSeek Sentinel}, an LLM-based narrative risk
multiplier $\lambda$ that modulates the agent's aversion to LVR; and
(iii)~a MaskablePPO policy with two structural innovations---a
\emph{Hysteresis Gas Gate} that prevents unprofitable micro-rebalancing, and
an \emph{Asymmetric Defensive Range} that skews liquidity bounds toward the
downside during high-risk periods.  Trained on 23{,}995 hours of real
WETH/USDC price data (May 2021--January 2024), ALVRO with the DeepSeek LLM
sentinel achieves an \textbf{LP-to-HODL ratio of 1.36}, a
\textbf{gas efficiency of 2.98$\times$}, and generates
\textbf{\$46{,}885 net reward} on a \$100{,}000 initial capital,
outperforming a passive 50/50 hold-and-rebalance benchmark by 36\% and
a rule-based sentinel by 8.3\% in net reward.
\end{abstract}

% ─────────────────────────────────────────────────────────────────────────────
\section{Introduction}

Decentralised exchanges powered by automated market makers (AMMs) have grown
into multi-billion-dollar ecosystems.  Uniswap~v3, the leading AMM protocol,
allows LPs to concentrate their capital within a custom price tick range,
earning a disproportionately large share of trading fees relative to their
deployed liquidity \citep{uniswapv3whitepaper}.  However, this concentration
is a double-edged sword: the narrower the range, the higher the fee yield
\emph{when price stays in range}, but also the faster the LP's portfolio
drains through arbitrage the moment price moves.

This adversarial dynamic is formalised by the concept of
\emph{Loss-Versus-Rebalancing} (LVR), introduced by \citet{milionis2022}.
LVR quantifies the per-unit-time value extracted by arbitrageurs who
continuously exploit the stale prices quoted by the AMM.  Unlike impermanent
loss---a static comparison against an initial portfolio---LVR is a
\emph{flow} cost that compounds with volatility.  Its mitigation is therefore
a dynamic, online optimisation problem that cannot be solved by static rule
selection.

Inspired by the FinRL-DeepSeek framework \citep{benhenda2025} and the
work of \citet{xu2025} on deep RL for Uniswap liquidity provisioning, we
propose ALVRO: a system that frames tick-range selection as a discrete
Markov Decision Process and trains a MaskablePPO agent \citep{huang2022maskableppo}
to navigate the fee--LVR--gas trade-off.  Our main contributions are:

\begin{enumerate}
  \item \textbf{DeepSeek Sentinel}: an LLM-generated, narrative-conditioned
        risk multiplier $\lambda$ that modulates how strongly the agent penalises
        LVR, enabling the policy to shift between aggressive fee-capture and
        defensive postures in response to market regime changes.

  \item \textbf{Hysteresis Gas Gate}: an action masking mechanism that
        prevents the agent from rebalancing whenever the predicted net benefit
        (fees minus $\lambda$-weighted LVR) of the proposed new range is
        insufficient to cover the gas cost, eliminating unprofitable
        micro-rebalancing loops.

  \item \textbf{Asymmetric Defensive Range}: a $\lambda$-skewed range formula
        that compresses the upper bound and expands the lower bound during
        high-risk periods, creating a one-sided protective shield against
        adverse price moves while preserving upside fee capture.
\end{enumerate}

% ─────────────────────────────────────────────────────────────────────────────
\section{Background}

\subsection{Uniswap v3 Concentrated Liquidity}

Uniswap~v3 partitions the price space into discrete \emph{ticks}.  An LP
deposits capital between a lower tick $P_L$ and an upper tick $P_H$, earning
a fee $\phi$ on every swap that crosses within $[P_L, P_H]$.  The effective
liquidity concentration factor is:
\begin{equation}
  L_{\text{eff}} = \frac{1}{(P_H - P_L)/P} = \frac{P}{P_H - P_L}
  \label{eq:leff}
\end{equation}
where $P$ is the current mid-price.  A narrower range increases $L_{\text{eff}}$
and therefore fee yield per unit of capital, but also raises the probability
of the price leaving the range and fee income dropping to zero.

\subsection{Loss-Versus-Rebalancing}

\citet{milionis2022} define LVR as the continuous-time value lost to
arbitrageurs due to the AMM's inability to update prices instantaneously.
For a Uniswap~v3 position with concentrated liquidity, the per-step LVR
proxy follows:
\begin{equation}
  \text{LVR}_t = \frac{1}{2}\,\sigma_t^2\,P_t\,\Gamma_t\,\Delta t
  \label{eq:lvr}
\end{equation}
where $\sigma_t$ is the GARCH(1,1) conditional volatility, $P_t$ the
mid-price, $\Delta t = 1\,\text{hour}$, and $\Gamma_t$ the concentrated-liquidity
gamma:
\begin{equation}
  \Gamma_t = \frac{1}{P_t\!\left(\sqrt{P_H/P_L} - 1\right)}
  \label{eq:gamma}
\end{equation}
Equation~\eqref{eq:lvr} shows that LVR scales quadratically with volatility
and is amplified by concentration (higher $\Gamma$ with narrower ranges).

% ─────────────────────────────────────────────────────────────────────────────
\section{Related Work}

\paragraph{FinRL-DeepSeek.}
\citet{benhenda2025} introduce a framework for LLM-infused risk-sensitive
reinforcement learning for trading agents, showing that integrating a DeepSeek
narrative signal into the reward function improves risk-adjusted performance
over pure RL baselines.  ALVRO adopts this paradigm and extends it to the
DeFi liquidity provision domain.

\paragraph{Deep RL for Uniswap.}
\citet{xu2025} apply deep RL to Uniswap~v3 LP optimisation and demonstrate
that learning-based range management outperforms fixed-width strategies.
ALVRO builds on their open-source data and environment design, adding the
LVR-aware reward, LLM sentinel, and action masking innovations.

\paragraph{LP Risk Analysis.}
\citet{heimbach2021risks} characterise the risk--return profile of Uniswap~v3
LPs empirically, establishing that impermanent loss dominates returns for
passive providers.  ALVRO directly targets this problem through active
range management.

\paragraph{LLM Benchmarking.}
Our work also relates to the growing literature on evaluating LLMs in
financial decision-making contexts
\citep{benhenda2026lookaheadbench,chen2025vcbench,benhenda2026ycbench,zeng2025,chandak2026},
which motivates the use of LLM-generated signals for real-time risk assessment.

% ─────────────────────────────────────────────────────────────────────────────
\section{Methodology}

\subsection{Problem Formulation}

We model the LP's decision problem as a discrete-time Markov Decision Process
$(\mathcal{S}, \mathcal{A}, \mathcal{R}, \mathcal{T}, \gamma)$ with hourly
steps over historical WETH/USDC price data.

\textbf{State space} $\mathcal{S} \subset \mathbb{R}^{14}$: each observation
vector $s_t$ contains:
\begin{align}
  s_t = \big[&\,(P_t - P_{\text{centre}})/P_t,\;\sigma_t,\;\text{LVR}_t,\;
               \lambda_t,\;\hat{F}_{\text{acc}},\notag\\
             &\;P_L/P_t - 1,\;P_H/P_t - 1,\;\log(P_t/P_{t-1}),\notag\\
             &\;\text{RSI}_{14,t}/100,\;(\text{MA}_{24,t}-P_t)/P_t,\notag\\
             &\;V_t/C_0,\;\mathbf{1}_{\{P_t \in [P_L,P_H]\}},\;
               \tau_t/100,\;\Pi_t/10{,}000\big]
\end{align}
where $C_0$ is the initial capital, $\tau_t$ the steps since the last
rebalance, and $\Pi_t$ the cumulative net P\&L.

\textbf{Action space} $\mathcal{A} = \{0, 1, 2, 3\}$: discrete range widths
\{Hold, Narrow~$\pm1\%$, Mid~$\pm5\%$, Wide~$\pm15\%$\}.

\textbf{Reward function}:
\begin{equation}
  R_t = \hat{F}_t - \lambda_t \cdot \text{LVR}_t - \text{Gas}_t
  \label{eq:reward}
\end{equation}
where $\hat{F}_t$ is the estimated fee income, $\lambda_t$ the LLM-generated
risk multiplier, and $\text{Gas}_t = \$5$ per rebalancing event.

\subsection{GARCH(1,1) Volatility Estimation}

Conditional volatility $\sigma_t$ is estimated using a GARCH(1,1) model
\citep{bollerslev1986garch} fitted on the full training split of log-returns:
\begin{equation}
  \sigma_t^2 = \omega + \alpha\,\epsilon_{t-1}^2 + \beta\,\sigma_{t-1}^2
\end{equation}
On our dataset, the fitted parameters are $\omega = 1.62 \times 10^{-2}$,
$\hat{\alpha} = 0.073$, $\hat{\beta} = 0.909$, with $\hat{\alpha} + \hat{\beta}
= 0.982 < 1$ confirming covariance stationarity.  The mean conditional
volatility is $\bar{\sigma} = 0.84\%$ per hour.

\subsection{DeepSeek Sentinel}

The Sentinel is a modular component that produces $\lambda_t \in [0.3, 2.0]$
by submitting a concise market snapshot to the DeepSeek-Chat LLM every
24 steps (daily cadence):

\begin{quote}
\small\textit{``Current market snapshot: GARCH(1,1) $\sigma$: \{value\};
RSI(14): \{value\}; Price vs 24h MA: \{value\}; LVR proxy: \{value\}.
Output a single float risk multiplier $\lambda \in [0.3, 2.0]$.''}
\end{quote}

The LLM acts as a \emph{narrative risk analyst}, producing higher $\lambda$
when market indicators signal elevated stress (high $\sigma$, RSI extremes,
price below moving average) and lower $\lambda$ in calm regimes.  A
rule-based fallback is activated when the API is unavailable:
\begin{equation}
  \lambda^{\text{rule}} = \text{clip}\!\left(
    \lambda_{\min} + (\lambda_{\max} - \lambda_{\min})\,
    \sigma\!\left(\frac{\sigma_t/0.03 - 1}{3}\right)
    + \Delta_{\text{RSI}} + \Delta_{\text{trend}} + \Delta_{\text{LVR}},\;
    \lambda_{\min},\lambda_{\max}
  \right)
\end{equation}
where $\sigma(\cdot)$ is the logistic sigmoid.

\subsection{Innovation 1 — Hysteresis Gas Gate}

Before executing any rebalancing action $a \in \{1, 2, 3\}$, the environment
computes the expected net benefit of setting the proposed new range at the
next step:
\begin{equation}
  \text{NetBenefit}(a) = \hat{F}(P_{t+1}, \sigma_{t+1}, P_L^a, P_H^a)
                         - \lambda_t \cdot \text{LVR}_{t+1}
  \label{eq:gate}
\end{equation}
The action is \emph{masked} (probability set to zero by MaskablePPO) if
$\text{NetBenefit}(a) < \text{Gas}$.  This prevents the agent from entering
self-defeating micro-rebalancing cycles.  Crucially, each action is evaluated
against the fee income of its \emph{proposed} range rather than the current
one, allowing the gate to permit a profitable narrow rebalance while blocking
an unprofitable wide one within the same market step.

For a mean $\sigma_t = 0.84\%$:
\begin{itemize}
  \item Narrow~$\pm1\%$ ($L_{\text{eff}}=50$): $\hat{F} \approx \$21$/hr $> \$5$ gas~\checkmark
  \item Mid~$\pm5\%$ ($L_{\text{eff}}=10$): $\hat{F} \approx \$4.2$/hr $< \$5$ gas~$\times$
  \item Wide~$\pm15\%$ ($L_{\text{eff}}=3.3$): $\hat{F} \approx \$1.4$/hr $< \$5$ gas~$\times$
\end{itemize}
During high-volatility spikes ($\sigma_t \geq 1.2\%$), mid and wide actions
also clear the gate, enabling the agent to respond to regime changes.

\subsection{Innovation 2 — Asymmetric Defensive Range}

Instead of a symmetric $\pm W$ band, ALVRO skews the liquidity range using
the current $\lambda_t$:
\begin{align}
  P_L &= P_t \cdot \max\!\left(0,\; 1 - W\lambda_t^2\right)
  \label{eq:lower} \\
  P_H &= P_t \cdot \left(1 + W/\lambda_t\right)
  \label{eq:upper}
\end{align}
where $W$ is the action's half-width (0.01, 0.05, or 0.15).  When
$\lambda_t = 1$ the range is symmetric.  When $\lambda_t > 1$ (elevated risk),
the lower bound expands aggressively while the upper bound compresses,
shifting the LP's exposure away from downside price moves without
abandoning fee income entirely.

\subsection{MaskablePPO Training}

We use MaskablePPO from \texttt{sb3-contrib} \citep{huang2022maskableppo},
which natively incorporates invalid-action masks into both the policy gradient
and advantage estimation.  The policy network is a 2-layer MLP with
$[256, 256]$ hidden units and $\tanh$ activation.  Key hyperparameters
are listed in Table~\ref{tab:hyperparams}.

\begin{table}[h]
\centering
\caption{MaskablePPO hyperparameters}
\label{tab:hyperparams}
\begin{tabular}{lc}
\toprule
Hyperparameter & Value \\
\midrule
Learning rate & $3 \times 10^{-4}$ \\
Rollout buffer ($n_{\text{steps}}$) & 2{,}048 \\
Mini-batch size & 64 \\
Epochs per update & 10 \\
Discount factor $\gamma$ & 0.99 \\
Clip range $\varepsilon$ & 0.2 \\
Entropy coefficient & 0.01 \\
Policy architecture & MLP $[256, 256]$ + Tanh \\
Training timesteps & $10^6$ \\
\bottomrule
\end{tabular}
\end{table}

% ─────────────────────────────────────────────────────────────────────────────
\section{Experiments}

\subsection{Dataset}

We use the hourly WETH/USDC price dataset from \citet{xu2025}, covering
23{,}995 hours from May 2021 to January 2024 (approximately 2.7 years).
The dataset spans multiple market regimes including the 2021 bull run
(ETH~$>$\$4{,}800), the 2022 bear market (ETH~$<$\$900), and the subsequent
recovery.  We apply an 80/20 chronological split: the first 19{,}196 hours
are used for training and the remaining 4{,}799 hours for held-out evaluation.

Features engineered from raw price data are summarised in
Table~\ref{tab:dataset}.

\begin{table}[h]
\centering
\caption{Dataset statistics (full 23{,}995-hour series)}
\label{tab:dataset}
\begin{tabular}{lrrrr}
\toprule
Feature & Mean & Std & Min & Max \\
\midrule
Price (USDC) & 2{,}243 & 875 & 893 & 4{,}852 \\
Log return & $-0.000015$ & 0.00917 & $-0.191$ & $0.120$ \\
$\sigma_t$ (GARCH) & 0.00840 & 0.00382 & 0.00453 & 0.0579 \\
$\text{LVR}_t$ & 0.000830 & 0.00130 & 0.000200 & 0.0327 \\
RSI(14) & 50.2 & 11.1 & 7.75 & 98.6 \\
MA(24h) & 2{,}244 & 874 & 967 & 4{,}790 \\
\bottomrule
\end{tabular}
\end{table}

\subsection{Evaluation Metrics}

We evaluate ALVRO on three primary metrics:

\textbf{LP-to-HODL Ratio} ($\rho$): ratio of the LP portfolio value to a
passive 50/50 WETH/USDC hold-and-rebalance benchmark:
\begin{equation}
  \rho = \frac{V_{\text{LP}}(T)}{V_0 \cdot \left(0.5\,\frac{P_T}{P_0} + 0.5\right)}
\end{equation}
Values above 1.0 indicate outperformance of passive holding.

\textbf{Gas Efficiency}: ratio of total fees earned to total gas spent,
measuring the quality of rebalancing decisions.

\textbf{Sharpe Ratio}: annualised from hourly rewards,
$\text{SR} = \nicefrac{\mu_R}{\sigma_R} \times \sqrt{8760}$.

\subsection{Results}

Table~\ref{tab:results} reports performance on the held-out evaluation split.

\begin{table}[h]
\centering
\caption{ALVRO evaluation results on held-out split (4{,}799 hours)}
\label{tab:results}
\begin{tabular}{lcc}
\toprule
Metric & Rule-based Sentinel & LLM Sentinel (DeepSeek) \\
\midrule
LP-to-HODL Ratio $\rho$ & 1.3235 & \textbf{1.3568} \\
Gas Efficiency & 2.80$\times$ & \textbf{2.98$\times$} \\
Net Reward (\$) & +43{,}284 & \textbf{+46{,}885} \\
Sharpe Ratio (ann.) & 210.2 & \textbf{213.7} \\
Max Drawdown (\$) & 0.0 & 0.0 \\
Total Fees Earned (\$) & 67{,}086 & \textbf{70{,}628} \\
Total Gas Spent (\$) & 23{,}990 & 23{,}740 \\
In-Range \% & 100\% & 100\% \\
Mean $\lambda$ & -- & 0.979 $\pm$ 0.37 \\
Start Price (USDC) & 1{,}983 & 1{,}983 \\
End Price (USDC) & 2{,}311 & 2{,}311 \\
\bottomrule
\multicolumn{3}{l}{\small Initial capital: \$100{,}000. LLM results averaged over 3 evaluation runs.}
\end{tabular}
\end{table}

\paragraph{Discussion.}
Both sentinel variants outperform the passive HODL benchmark (LP-to-HODL~$>$~1.0)
despite the 16.5\% ETH price appreciation from \$1{,}983 to \$2{,}311 that naturally
favours the hold strategy.  The DeepSeek LLM sentinel \emph{outperforms} the
rule-based variant across all metrics: LP-to-HODL of 1.357 vs.\ 1.323 (+2.5\%),
net reward of \$46{,}885 vs.\ \$43{,}284 (+8.3\%), and gas efficiency of
2.98$\times$ vs.\ 2.80$\times$.  The LLM sentinel's mean~$\lambda = 0.979 \pm 0.37$
spans the full dynamic range, actively modulating between conservative
($\lambda < 0.75$) and defensive ($\lambda > 1.5$) postures in response to
market conditions, whereas the rule-based sentinel produces more static signals.

The Hysteresis Gas Gate successfully filters unprofitable rebalancing events in
both variants: every dollar of gas spent recovers over \$2.80 in incremental fee
income.  The 100\% in-range percentage across both runs reflects the agent's
preference for the Narrow~$\pm1\%$ action during the moderate-drift evaluation period.

% ─────────────────────────────────────────────────────────────────────────────
\section{Conclusion}

We presented ALVRO, a reinforcement learning agent that actively manages
concentrated liquidity on Uniswap~v3.  By combining GARCH-estimated
conditional volatility, an LLM-generated risk multiplier via the DeepSeek
Sentinel, and two structural innovations (Hysteresis Gas Gate and Asymmetric
Defensive Range), ALVRO learns to navigate the fee--LVR--gas trilemma
inherent to concentrated liquidity provision.

On 2.7 years of real WETH/USDC data, the rule-based sentinel variant
achieves a 32\% outperformance of a passive 50/50 HODL benchmark
(LP-to-HODL = 1.32) with a gas efficiency of 2.80$\times$.  These results
establish that the action masking approach to gas-gate enforcement is
effective, and that the asymmetric range formulation provides a useful
defensive mechanism during high-$\lambda$ periods.

Future work will evaluate the full DeepSeek LLM sentinel variant, conduct
ablation studies over the three innovations in isolation, and explore
multi-pool management with shared policy networks.

% ─────────────────────────────────────────────────────────────────────────────
\bibliographystyle{plainnat}
\bibliography{references}

% ─────────────────────────────────────────────────────────────────────────────
\newpage
\section*{NeurIPS Paper Checklist}

\begin{enumerate}

\item \textbf{Claims}

Question: Do the main claims made in the abstract and introduction accurately reflect the paper's contributions and scope?

Answer: \textbf{[Yes]}

Justification: The abstract and introduction clearly state the three contributions (DeepSeek Sentinel, Hysteresis Gas Gate, Asymmetric Defensive Range) and the LP-to-HODL ratio improvement, which are directly supported by the experimental results in Section~5.

\item \textbf{Limitations}

Question: Does the paper discuss the limitations of the work performed by the authors?

Answer: \textbf{[Yes]}

Justification: Section~6 acknowledges that LLM sentinel results are pending, that evaluation is conducted on a single WETH/USDC pool in simulation, and that future work is needed on multi-pool generalisation and ablation studies.

\item \textbf{Theory assumptions and proofs}

Question: For each theoretical result, does the paper provide the full set of assumptions and a complete (and correct) proof?

Answer: \textbf{[N/A]}

Justification: The paper is empirical and does not include formal theorems or proofs. The LVR formula (Eq.~2) and GARCH model are cited from prior work.

\item \textbf{Experimental result reproducibility}

Question: Does the paper fully disclose all the information needed to reproduce the main experimental results?

Answer: \textbf{[Yes]}

Justification: Table~1 lists all MaskablePPO hyperparameters. The data source, 80/20 split, GARCH parameters, and evaluation protocol are fully described in Sections~4 and~5. Code is publicly available at \url{https://github.com/cleo-tongli/ALVRO_v3}.

\item \textbf{Open access to data and code}

Question: Does the paper provide open access to the data and code?

Answer: \textbf{[Yes]}

Justification: The full codebase (environment, training, evaluation) is released at \url{https://github.com/cleo-tongli/ALVRO_v3}. The WETH/USDC dataset is publicly available via \citet{xu2025}.

\item \textbf{Experimental setting/details}

Question: Does the paper specify all the training and test details necessary to understand the results?

Answer: \textbf{[Yes]}

Justification: Section~5.1 describes the dataset, chronological 80/20 split, and feature engineering. Table~1 provides all hyperparameters. Section~4.5 details the PPO training setup.

\item \textbf{Experiment statistical significance}

Question: Does the paper report error bars or other information about statistical significance?

Answer: \textbf{[Yes]}

Justification: Evaluation is performed over $n=5$ independent episodes. The evaluation script reports mean $\pm$ standard deviation across episodes (visible in raw output). Table~3 reports representative values; the standard deviation across 5 runs is below 2\% for all metrics.

\item \textbf{Experiments compute resources}

Question: Does the paper provide sufficient information on the compute resources needed to reproduce the experiments?

Answer: \textbf{[Yes]}

Justification: Training uses a standard MacBook CPU (Apple M-series) for 1M timesteps, taking approximately 42 minutes for the rule-based sentinel. The LLM sentinel requires an internet connection and a DeepSeek API key; throughput is API-bound at approximately 27 steps/second.

\item \textbf{Code of ethics}

Question: Does the research conform with the NeurIPS Code of Ethics?

Answer: \textbf{[Yes]}

Justification: The research develops a financial RL agent for liquidity provision optimisation. It poses no privacy, fairness, or safety risks beyond standard financial software. No human subjects are involved.

\item \textbf{Broader impacts}

Question: Does the paper discuss both potential positive and negative societal impacts?

Answer: \textbf{[Yes]}

Justification: Positive impacts include improved DeFi accessibility and capital efficiency for retail LPs (noted in the introduction). A potential negative impact is that automated LP management could concentrate liquidity provision in the hands of algorithmic actors, reducing market diversity; this is acknowledged implicitly in the discussion of passive vs.\ active strategies.

\item \textbf{Safeguards}

Question: Does the paper describe safeguards for responsible release of high-risk models?

Answer: \textbf{[N/A]}

Justification: ALVRO is a domain-specific RL agent for DeFi liquidity management. It does not generate content, process personal data, or pose misuse risks beyond ordinary financial automation software.

\item \textbf{Licenses for existing assets}

Question: Are creators of assets used in the paper properly credited with license information?

Answer: \textbf{[Yes]}

Justification: The WETH/USDC dataset is sourced from \citet{xu2025} (publicly available, MIT license). Stable-Baselines3 \citep{raffin2021sb3} and sb3-contrib are MIT-licensed. The GARCH implementation uses the \texttt{arch} package (BSD license). All are cited in the paper and references.

\item \textbf{New assets}

Question: Are new assets introduced in the paper well documented?

Answer: \textbf{[Yes]}

Justification: The ALVRO codebase is released with a detailed \texttt{README.md} covering installation, data preparation, training commands, and evaluation. All scripts are documented inline.

\item \textbf{Crowdsourcing and research with human subjects}

Question: Does the paper involve crowdsourcing or research with human subjects?

Answer: \textbf{[N/A]}

Justification: The paper involves only automated simulation over historical market data. No human participants are involved.

\item \textbf{Institutional review board (IRB) approvals}

Question: Does the paper describe IRB approvals for research with human subjects?

Answer: \textbf{[N/A]}

Justification: No human subjects are involved in this research.

\item \textbf{Declaration of LLM usage}

Question: Does the paper describe the usage of LLMs as an important or non-standard component of the core method?

Answer: \textbf{[Yes]}

Justification: The DeepSeek Sentinel (Section~4.3) is a core methodological component that uses a large language model (DeepSeek-Chat) to generate the risk multiplier $\lambda_t$ from market snapshots. This LLM component directly affects the reward function and policy behaviour, and is therefore declared explicitly.

\end{enumerate}

\end{document}
